\SixthArticle{22}
Die Funktion $\Pi(k,z)$ mag zwar allgemeiner erscheinen als $\Pi z$, wird uns aber von nun an entbehrlich sein, da sie sich leicht auf die letztere zurückführen lässt. Aus der Verbindung der Gleichungen 38, 45, 46 ergibt sich nämlich
\[
\begin{array}{@{}ll@{}}
[47]&\displaystyle\Pi(k,z)=\frac{k^z\Pi k.\Pi z}{\Pi(k+z)}
\end{array}
\]
Im Übrigen ist der Zusammenhang dieser Funktionen mit denen, die der berühmte \textsc{Kramp} \emph{numerische\SourceLeaf{156} Fakultäten} genannt hat, an sich offensichtlich. Die numerische Fakultät, die dieser Autor mit
$a^{bIc}$
bezeichnet, ist nämlich in unseren Zeichen


\[
=\frac{c^b b^{\frac{a}{c}-1}\Pi b}{\Pi(b,\frac{a}{c}-1)}
=\frac{c^b\Pi(\frac{a}{c}+b-1)}{\Pi(\frac{a}{c}-1)}
\]
Es scheint aber ratsamer, eine Funktion \emph{einer} Veränderlichen in die Analysis einzuführen als eine Funktion dreier Veränderlicher, besonders da sich diese auf jene zurückführen lässt.

\SixthArticle{23}
Die Stetigkeit der Funktion $\Pi z$ wird unterbrochen, sooft ihr Wert unendlich groß wird, das heißt für negative ganzzahlige Werte von $z$. Sie wird daher von $z=-1$ bis $z=\infty$ positiv sein; da $\Pi z$ an beiden Grenzen einen unendlich großen Wert annimmt, wird es dazwischen einen kleinsten Wert geben, den wir als $=0{,}8856024$ gefunden haben und der dem Wert $z=0{,}4616321$ entspricht. Zwischen den Grenzen $z=-1$ und $z=-2$ wird der Wert der Funktion $\Pi z$ negativ, zwischen $z=-2$ und $z=-3$ wieder positiv und so fort, wie sich unmittelbar aus Gleichung 44 ergibt. Ferner ist klar: Darf man sämtliche Werte der Funktion $\Pi z$ zwischen beliebigen, um die Einheit verschiedenen Grenzen, zum Beispiel von $z=0$ bis $z=1$, als bekannt ansehen, so lässt sich daraus mit Hilfe der Gleichung 45 der Wert der Funktion für jeden anderen reellen Wert von $z$ leicht herleiten. Zu diesem Zweck haben wir eine \emph{Tafel} angefertigt, die am Ende dieses Abschnitts angefügt ist und die Briggsschen Logarithmen der Funktion $\Pi z$ für $z=0$ bis $z=1$ in Schritten von einem Hundertstel mit zwanzig Stellen angibt, mit größter Sorgfalt berechnet. Dabei ist jedoch zu bemerken, dass die letzte, zwanzigste Stelle bisweilen um eine oder zwei Einheiten fehlerhaft sein kann.

